% Main-text-oriented matrix reformulation of the current OODKA-OT code.
% Content follows docs/oodka_current_method.tex; implementation-specific
% constants are retained in the appendix rather than interrupting the method.

\documentclass[10pt]{article}

\usepackage[margin=1in]{geometry}
\usepackage[T1]{fontenc}
\usepackage[utf8]{inputenc}
\usepackage{amsmath,amssymb,mathtools}
\usepackage{bm}
\usepackage{booktabs}
\usepackage{graphicx}
\usepackage{microtype}
\usepackage[hidelinks]{hyperref}

\newcommand{\R}{\mathbb{R}}
\newcommand{\cI}{\mathcal{I}}
\newcommand{\cQ}{\mathcal{Q}}
\newcommand{\cL}{\mathcal{L}}
\newcommand{\cG}{\mathcal{G}}
\newcommand{\ones}{\bm{1}}
\newcommand{\diag}{\operatorname{Diag}}
\newcommand{\Mat}{\operatorname{Mat}}
\newcommand{\Resize}{\operatorname{Resize}}
\newcommand{\IN}{\operatorname{IN}}
\newcommand{\LN}{\operatorname{LN}}
\newcommand{\GELU}{\operatorname{GELU}}
\newcommand{\softplus}{\operatorname{softplus}}
\newcommand{\sigm}{\operatorname{sigmoid}}
\newcommand{\Norm}{\operatorname{Norm}}
\newcommand{\clip}{\operatorname{clip}}
\newcommand{\KL}{\operatorname{KL}}
\newcommand{\stopgrad}{\operatorname{stopgrad}}
\newcommand{\rowcos}{\operatorname{rowcos}}
\newcommand{\RoiBox}{\operatorname{Box}}
\newcommand{\Crop}{\operatorname{CropResize}}
\newcommand{\ind}{\mathbb{1}}
\newcommand{\balpha}{\bm{\alpha}}
\newcommand{\bbeta}{\bm{\beta}}

\title{OODKA-OT: Unified Matrix Formulation}
\date{}

\begin{document}
\maketitle

\begin{abstract}
OODKA-OT learns a prompt-indexed mixture of two complementary feature
components from a frozen BiomedParse foundation model.  A frozen nnUNet expert
is used only during training to define component-specific transport teachers:
balanced transport for the structure-oriented component $P$ and unbalanced
transport for the task-residual component $S$.  We formulate decomposition,
decorrelation, prompt-conditioned Beta routing, transport distillation, and LGE
ROI refinement as one factorize--align--recombine system.  The formulation is
matrix-based but remains faithful to the current implementation, including
valid-slice masking, loss reductions, detached transport, cached training ROIs,
and pure-student deployment.
\end{abstract}

\section{Unified Problem Formulation}

Let $B_{\phi_b}$ denote the frozen BiomedParse visual encoder and
$E_{\phi_n}$ the frozen nnUNet expert; the frozen BiomedParse pixel decoder and
prompt predictor are denoted separately below.  Their parameters
$\phi_b,\phi_n$ never receive gradients.  The trainable set $\Theta$ contains
the foundation/expert feature adapters, the two-component factorization
modules, and the prompt-conditioned router.

Training uses block tuples
\begin{equation}
 (\bm x^b,\bm x^n,\bm Y,\bm v)
 \in
 \R^{B\times Z\times3\times H\times W}
 \times\R^{B\times C_n\times Z\times H\times W}
 \times\{-1,0,1,\ldots\}^{B\times Z\times H\times W}
 \times\{0,1\}^{B\times Z},
 \label{eq:block-inputs}
\end{equation}
where $\bm v$ marks real versus repeated tail slices.  Both two-dimensional
encoders process the $BZ$ slices as a flattened batch before their features are
restored to block form.

We use $d\in\{b,n\}$ for the foundation/expert side,
$\xi\in\{p,s\}$ for the two components,
$\ell\in\cI=\{2,3,4,5\}$ for feature level,
$r\in\cQ=\{\mathrm{mask},5,4,3\}$ for frozen pixel-decoder resolution,
$c$ for prompt, and $j\in\{A,R\}$ for the LGE anatomy/refinement pass.

For block $b$ with $Z$ slices, channels are placed in columns:
\begin{equation}
 \bm X_{\ell b}^{d}
 =\Mat(\bm Z_{\ell b}^{d})
 \in\R^{N_\ell\times C_\ell^d},
 \qquad N_\ell=ZH_\ell W_\ell .
 \label{eq:token-convention}
\end{equation}
Before unfolding, the raw expert feature is trilinearly resized to
$(Z,H_\ell,W_\ell)$.  A diagonal matrix
$\bm V_{\ell b}$ repeats $v_{bz}\in\{0,1\}$ over spatial rows and removes padded
tail slices from segmentation, reconstruction, decorrelation, and OT losses.
It does not remove them from the forward pass; the router KL has no slice
factor and is evaluated once per prompt field.
All instance-normalization operators below act independently per block and
channel over that block's $ZHW$ sites; no normalization crosses blocks.

Let $\mathcal D_{\Theta,\xi}^b$ denote the foundation factorizer,
$\Psi_{\phi_b}$ the frozen pixel decoder, $\mathcal M_{\bm\Gamma}$ the
prompt--spatial mixer, $\mathcal H_{\phi_b}$ the frozen prompt predictor,
$\mathcal A_{Q_m}$ mask-query averaging, and $\mathcal U$ output-logit resizing.
The deployed student is the single operator
\begin{equation}
\boxed{
\begin{aligned}
 \bm F_\xi
 &=\Psi_{\phi_b}\!\left(
    \mathcal D_{\Theta,\xi}^b(B_{\phi_b}(\bm x^b))\right),
 \qquad \xi\in\{p,s\},\\
 \bm L_\Theta(\bm x^b,\bm T)
 &=\mathcal U\circ\mathcal A_{Q_m}\circ\mathcal H_{\phi_b}
   \left(
   \mathcal M_{\bar{\bm\Gamma}_\Theta(\bm T)}[\bm F_p,\bm F_s],
   \bm T\right).
\end{aligned}}
\label{eq:student-master}
\end{equation}
Here $\bm T\in\R^{C_{\mathrm{pr}}\times d_t}$ stacks frozen prompt embeddings.
The bar denotes the deterministic Beta mean used for validation and deployment;
during training it is replaced by a reparameterized sample from the same
router field.  Both cases are formalized in
Section~\ref{sec:matrix-routing}.
The expert and all OT/UOT operators are absent from
Eq.~\eqref{eq:student-master}; they only constrain its representation during
training.

At training time the shared factors induce two parallel, rather than serial,
graphs:
\begin{equation}
\boxed{
\begin{aligned}
 \mathcal B_{\mathrm{pred}}:\quad
 &(\bm P^b,\bm S^b)
 \xrightarrow{\ \Psi_{\phi_b}\ }
 (\bm F_p,\bm F_s)
 \xrightarrow{\ \bm\Gamma_\Theta(\bm T)\ }
 \bm L
 \longrightarrow(\cL_{\mathrm{seg}},\cL_{\mathrm{route}}),\\
 \mathcal B_{\mathrm{align}}:\quad
 &(\bm P^b,\bm S^b;\bm P^n,\bm S^n)
 \longrightarrow
 \left\{
 \begin{aligned}
  &(\cL_{\mathrm{ae}},\cL_{\mathrm{ort}}),\\[-1mm]
  &(\bm\Pi^p,\bm\Pi^s)
   \longrightarrow(\widetilde{\bm Y}^p,\widetilde{\bm Y}^s)
   \longrightarrow(\cL_{P\text{-OT}},\cL_{S\text{-UOT}}).
 \end{aligned}
 \right.
\end{aligned}}
\label{eq:parallel-training-graph}
\end{equation}
The alignment branch terminates in regularizers; no transported teacher is an
input to the predictive branch.

\section{Complementary Matrix Factorization}

\subsection{Two-sided factorization}

All channel maps in this subsection are bias-free $1\times1\times1$
convolutions, hence right matrix multiplications under
Eq.~\eqref{eq:token-convention}.  The foundation decomposition is
\begin{equation}
\boxed{
 \begin{bmatrix}\bm P_{\ell b}^b&\bm S_{\ell b}^b\end{bmatrix}
 =
 \bm X_{\ell b}^b
 \begin{bmatrix}\bm W_{\ell p}^b&\bm W_{\ell s}^b\end{bmatrix}.}
\label{eq:foundation-factor}
\end{equation}
Both outputs retain $C_\ell^b$ channels and receive no activation or
normalization.

For the expert, define
\begin{equation}
\boxed{
\begin{aligned}
 \bm U_{\ell b}
 &=\GELU\!\left(
   \IN_{\ell}^{\mathrm{blk}}
   (\bm X_{\ell b}^n\bm W_\ell^0)\right),\\
 \begin{bmatrix}\bm P_{\ell b}^n&\bm S_{\ell b}^n\end{bmatrix}
 &=
 \begin{bmatrix}
  \mathcal N_{\ell p}(\bm U_{\ell b}\bm W_{\ell p}^n)&
  \mathcal N_{\ell s}(\bm U_{\ell b}\bm W_{\ell s}^n)
 \end{bmatrix}.
\end{aligned}}
\label{eq:expert-factor}
\end{equation}
$\mathcal N_{\ell p},\mathcal N_{\ell s}$ are affine instance normalization at
levels $2$--$4$ and identities at level $5$; the shared instance normalization
in the first line remains active at every level.  The branch projections map
to $C_\ell^b$ channels.  These equations create two components but do not
guarantee their meanings: balanced $P$-OT and residual $S$-UOT assign the
operational roles below.

\subsection{Completeness and non-redundancy}

Stack the $B$ block matrices vertically and their row masks block-diagonally,
for example
$\bm X_\ell^d=\operatorname{col}_{b=1}^B(\bm X_{\ell b}^d)$ and
$\bm V_\ell=\operatorname{bdiag}_{b=1}^B(\bm V_{\ell b})$.  For any such
stacked feature matrix with $C$ columns, define the valid relative energy
\begin{equation}
 \mathfrak R(\bm A,\bm X^\star;\bm V)
 =
 \frac{\|\bm V(\bm A-\bm X^\star)\|_F^2/n_V}
      {\|\bm V\bm X^\star\|_F^2/n_V+10^{-8}},
 \qquad
 n_V=C\max(\operatorname{tr}\bm V,HW)
 =CHW\max\!\left(\sum_{b,z} v_{bz},1\right).
 \label{eq:relative-energy}
\end{equation}
Let
\[
 \widehat{\bm X}_{\ell}^b=\bm P_{\ell}^b+\bm S_{\ell}^b,\qquad
 \widehat{\bm X}_{\ell}^n
 =\mathcal R_{\ell p}(\bm P_{\ell}^n)
  +\mathcal R_{\ell s}(\bm S_{\ell}^n),
\]
where each expert decoder $\mathcal R_{\ell\xi}$ is
Conv--IN--GELU--Conv.  Applying the frozen pixel decoder to the complete
level collections gives
\[
 \bm F_p^r=\Psi_{\phi_b}^r(\{\bm P_\ell^b\}_{\ell\in\cI}),\quad
 \bm F_s^r=\Psi_{\phi_b}^r(\{\bm S_\ell^b\}_{\ell\in\cI}),\quad
 \bm F_{p+s}^r=\Psi_{\phi_b}^r(
   \{\bm P_\ell^b+\bm S_\ell^b\}_{\ell\in\cI}).
\]
These decoder outputs are stacked over blocks in the same way, and $\bm V_r$
denotes the valid-slice flag broadcast over the spatial rows at decoder
resolution $r$.  The implemented reconstruction functional is
\begin{equation}
\boxed{
 \cL_{\mathrm{ae}}
 =\frac{1}{9}\left[
   \sum_{\ell\in\cI}\sum_{d\in\{b,n\}}
   \mathfrak R(\widehat{\bm X}_{\ell}^d,\bm X_{\ell}^d;\bm V_{\ell})
   +\frac{1}{|\cQ|}\sum_{r\in\cQ}
   \mathfrak R(\bm F_p^r+\bm F_s^r,\bm F_{p+s}^r;\bm V_r)
 \right].}
\label{eq:complete-reconstruction}
\end{equation}
This preserves the code's unusual but exact reduction: eight level/model terms
plus one term already averaged across four pixel-decoder outputs, all divided
by nine.

Let $\operatorname{Std}_{\bm V}(\cdot)$ standardize every channel over valid
spatial rows using the implementation's sample-standard-deviation path when
all slices are valid and masked population-RMS path otherwise.  The full
$P$--$S$ cross-correlation matrix is
\[
 \bm C_{\ell b}^{d}
 =\frac{1}{N_{\ell b}^{\mathrm{valid}}}
   \operatorname{Std}_{\bm V}(\bm P_{\ell b}^{d})^\top
   \operatorname{Std}_{\bm V}(\bm S_{\ell b}^{d}).
\]
Hence the separation penalty is the matrix norm
\begin{equation}
\boxed{
 \cL_{\mathrm{ort}}
 =\sum_{\ell\in\cI}\sum_{d\in\{b,n\}}
  \frac{1}{B(C_\ell^b)^2}
  \sum_{b=1}^{B}\|\bm C_{\ell b}^{d}\|_{1,1}.}
\label{eq:matrix-correlation}
\end{equation}
All $C_\ell^2$ cross-component channel pairs are penalized in magnitude.  This
is decorrelation, not strict subspace orthogonality or statistical independence.

\section{Prompt-conditioned Stochastic Matrix Routing}
\label{sec:matrix-routing}

Let
\[
 \bm H=\GELU(\LN(\bm T)\bm W_h+\ones\bm b_h^\top),\qquad
 \bm A^\alpha=\operatorname{Head}_\alpha(\bm H),\quad
 \bm A^\beta=\operatorname{Head}_\beta(\bm H).
\]
On the mask grid, define the RBF dictionary
$\bm\Phi\in\R^{N_0\times K}$ by
$\Phi_{ik}=\exp[-\|\bm u_i-\bm\mu_k\|_2^2/(2\sigma_{\mathrm{rbf}}^2)]$.
The complete prompt-by-position Beta field is generated by two matrix products:
\begin{equation}
\boxed{
\begin{aligned}
 \bm\eta^\alpha
 &=b_\alpha\ones_{C_{\mathrm{pr}}}\ones_{N_0}^\top
   +K^{-1/2}\bm A^\alpha\bm\Phi^\top,
 &\balpha&=\ones_{C_{\mathrm{pr}}}\ones_{N_0}^\top
   +\softplus(\bm\eta^\alpha),\\
 \bm\eta^\beta
 &=b_\beta\ones_{C_{\mathrm{pr}}}\ones_{N_0}^\top
   +K^{-1/2}\bm A^\beta\bm\Phi^\top,
 &\bbeta&=\ones_{C_{\mathrm{pr}}}\ones_{N_0}^\top
   +\softplus(\bm\eta^\beta).
\end{aligned}}
\label{eq:beta-matrix}
\end{equation}
The defaults are $G=8$, $K=G^2=64$ and
$\sigma_{\mathrm{rbf}}=3/(G-1)$.  Both coefficient heads, including their
biases, start at zero, while
$b_\alpha=\softplus^{-1}(6)$ and $b_\beta=\softplus^{-1}(2)$; consequently the
initial field is exactly $\operatorname{Beta}(7,3)$.

Training samples elementwise
$\bm\Gamma_b\sim\operatorname{Beta}(\balpha,\bbeta)$ independently across
blocks, prompts, and positions, and shares that realization across the $Z$
slices of block $b$.  Validation and deployment use
\begin{equation}
 \bar{\bm\Gamma}_\Theta(\bm T)
 =\mathbb{E}[\bm\Gamma]
 =\balpha\oslash(\balpha+\bbeta).
 \label{eq:beta-expectation}
\end{equation}
One mask-grid realization is area-resized to every $r\in\cQ$; coarse levels are
not independently resampled.

For
$\bm F_{p,bz}^r,\bm F_{s,bz}^r\in\R^{N_r\times C_r}$ and gate vector
$\bm\gamma_{bc}^r\in\R^{N_r}$, prompt-conditioned fusion is
\begin{equation}
\boxed{
 \widetilde{\bm F}_{bzc}^r
 =\bm F_{s,bz}^r
  +\diag(\bm\gamma_{bc}^r)
   (\bm F_{p,bz}^r-\bm F_{s,bz}^r).}
\label{eq:diagonal-fusion}
\end{equation}
The gate is scalar across channels but varies by prompt and position.  The
router is not image-conditioned: its distribution depends only on frozen prompt
semantics and fixed coordinates.

The prior regularizer is the elementwise Beta KL averaged over the
$C_{\mathrm{pr}}\times N_0$ field,
\begin{equation}
 \cL_{\mathrm{route}}
 =\frac{
  \langle\ones_{C_{\mathrm{pr}}}\ones_{N_0}^\top,\,
  \KL[\operatorname{Beta}(\balpha,\bbeta)
  \,\|\,\operatorname{Beta}(7,3)]\rangle_F}
 {C_{\mathrm{pr}}N_0}.
 \label{eq:route-matrix-kl}
\end{equation}
It contains no batch or slice factor.  The $BZC_{\mathrm{pr}}$ visual--prompt
pairs are flattened in $(b,z,c)$ order.  For predictor query tensor
$\bm Q_{bzc}\in\R^{Q_m\times N_o}$,
\[
 \bm L_{bzc}=Q_m^{-1}\ones_{Q_m}^\top\bm Q_{bzc},
\]
with no object-existence weighting, followed by trilinear logit resizing to
$(Z,H,W)$.

\section{Prompt-wise Segmentation as a Matrix Reduction}

For prompt $c$, let
$\bm y_{bc}=\operatorname{vec}(\ind[\bm Y_b=k_c])$ and let $\bm m_b$ mask
target $-1$ and invalid slices.  Denote masked BCE-with-logits and soft Dice by
$\mathfrak B_{bc}$ and $\mathfrak D_{bc}$.  Define
\begin{equation}
\begin{aligned}
 \Delta^{\mathrm{emp}}_{bc}
 &=\ind\!\left[
 \bm m_b^\top\bm y_{bc}
 <5\times10^{-4}\,\ones^\top\bm m_b\right],\\
 S_{bc}
 &=\mathfrak B_{bc}
 +(1-\Delta^{\mathrm{emp}}_{bc})\mathfrak D_{bc},\\
 W_{bc}
 &=\Omega_{bc}\left(1-\tfrac12\Delta^{\mathrm{emp}}_{bc}\right),
\end{aligned}
\label{eq:pair-matrix}
\end{equation}
where $\Omega_{bc}$ is prompt validity.  Thus an empty target uses BCE only and
retains weight $1/2$ rather than being discarded.

Collecting $S_{bc},W_{bc}$ into matrices
$\bm S,\bm W\in\R^{B\times C_{\mathrm{pr}}}$ gives
both executable reductions:
\begin{equation}
\boxed{
\begin{aligned}
 \cL_{\mathrm{seg}}^{\mathrm{mean}}
 &=\frac{\langle\bm W,\bm S\rangle_F}
 {\max(\langle\bm W,
 \ones_B\ones_{C_{\mathrm{pr}}}^\top\rangle_F,10^{-6})},\\
 \cL_{\mathrm{seg}}^{\mathrm{LGE}}
 &=\ones_{C_{\mathrm{pr}}}^\top\left[
  (\bm W\odot\bm S)^\top\ones_B
  \oslash\max(\bm W^\top\ones_B,
  10^{-6}\ones_{C_{\mathrm{pr}}})\right].
\end{aligned}}
\label{eq:seg-matrix-reduction}
\end{equation}
Generic training averages valid block--prompt pairs globally.  LGE averages
blocks separately for each prompt and then sums prompts; vector-valued maxima
and divisions in Eq.~\eqref{eq:seg-matrix-reduction} are elementwise.

\section{Component-aware Transport in One Variational System}

For valid slice $m$ at level $\ell$, adaptive average pooling gives
\begin{equation}
 \bm X_{\ell m}^\xi,\bm Y_{\ell m}^\xi
 \in\R^{N_\ell^{\mathrm{OT}}\times C_\ell},
 \qquad
 N_\ell^{\mathrm{OT}}
 =\min(H_\ell,32)\min(W_\ell,32).
 \label{eq:ot-tokens}
\end{equation}
$\bm X$ and $\bm Y$ are foundation and channel-aligned expert tokens.  Both
sides share the same grid, and independent slices never share a plan.

\subsection{Branch-specific masses}

For a nonnegative $N$-vector, write
\[
 \Norm(\bm q)=
 \begin{cases}
  \bm q/(\ones_N^\top\bm q),&\ones_N^\top\bm q>10^{-8},\\
  \ones_N/N,&\text{otherwise}.
 \end{cases}
\]
For class mask $\bm M_c=\ind[\bm Y=k_c]$, let
$\bm O,\bm B,\bm M_{\mathrm{res}}$ collect pooled occupancy, boundary, and
max-pooled rescue maps by class, and let
\[
 \bar\pi_c=
 \begin{cases}|M_c|^{-1/2},&|M_c|>0,\\0,&|M_c|=0,\end{cases}
 \qquad
 \bm\pi=\Norm(\bar{\bm\pi}),\qquad
 \bm h=(\bm O+1.5\bm B+0.2\bm M_{\mathrm{res}})\bm\pi .
\]
Define stabilized row magnitude, normalized energy, and $S$ specificity:
\[
 \bm\nu(\bm F)_i=\sqrt{\|\bm F_{i:}\|_2^2+\varepsilon_{\mathrm{num}}},
 \quad
 \mathcal E(\bm F)=
 \frac{\bm\nu(\bm F)}
 {\max(\operatorname{mean}\bm\nu(\bm F),\varepsilon_{\mathrm{num}})},
 \quad
 \bm\chi(\bm P,\bm S)=
 \frac{\bm\nu(\bm S)}
 {\bm\nu(\bm P)+\bm\nu(\bm S)+\varepsilon_{\mathrm{num}}\ones},
\]
where $\varepsilon_{\mathrm{num}}=10^{-8}$.  The structure marginals are
\begin{equation}
\boxed{
\begin{aligned}
 \bm a^p&=\Norm\!\left(
  \bm h\odot[\ones+0.1\mathcal E(\bm X^p)]\right),\\
 \bm b^p&=\Norm\!\left(
  \bm h\odot[\ones+0.1\mathcal E(\bm Y^p)]\right).
\end{aligned}}
\label{eq:p-mass-matrix}
\end{equation}
If $\ones^\top\bm h\leq10^{-8}$, each side falls back to its own energy; every
numerically zero candidate falls back to a uniform mass.

Detached prompt-valid BCE maps are averaged as
\[
 \bm d_{\mathrm{pix}}
 =\frac{\sum_c\Omega_c\bm e_c^b}{\max(\sum_c\Omega_c,1)},\qquad
 \bm e^n_{\mathrm{pix}}
 =\frac{\sum_c\Omega_c\bm e_c^n}{\max(\sum_c\Omega_c,1)}
\]
and pooled to $\bm d,\bm e^n$.  With
$\bm\Delta_{\mathrm{adv}}=\bm d-\bm e^n$,
\[
 \widehat{\bm d}
 =\clip\!\left(\bm d/\max(\operatorname{mean}\bm d,10^{-8}),0,3\right),
 \quad
 s_\Delta=\max(\operatorname{mean}|\bm\Delta_{\mathrm{adv}}|,10^{-8}),
 \quad
 \bm\zeta
 =\clip\!\left(2\sigm(\bm\Delta_{\mathrm{adv}}/(0.5s_\Delta)),0,3\right).
\]
The residual marginals are
\begin{equation}
\boxed{
\begin{aligned}
 \bm a^s&=\Norm\!\left(
  \mathcal E(\bm X^s)\odot\bm\chi(\bm X^p,\bm X^s)
  \odot(0.1\ones+\widehat{\bm d})\right),\\
 \bm b^s&=\Norm\!\left(
  \mathcal E(\bm Y^s)\odot\bm\chi(\bm Y^p,\bm Y^s)
  \odot(0.1\ones+\bm\zeta)\right).
\end{aligned}}
\label{eq:s-mass-matrix}
\end{equation}
For grouped LGE labels, expert class probabilities are summed within each
source group, clipped to $[10^{-6},1-10^{-6}]$, and converted to a binary logit
before its BCE map is formed.

\subsection{Cost, transport, and teacher}

Let $\operatorname{RowNorm}$ be detached stabilized row normalization and
$\bm U$ the normalized coordinate matrix.  The two costs are one family:
\begin{equation}
\boxed{
\begin{aligned}
 \bm C^\xi
 ={}&
 \left[\ones_{N_\ell^{\mathrm{OT}}}
 \ones_{N_\ell^{\mathrm{OT}}}^\top-
 \operatorname{RowNorm}(\bm X^\xi)
 \operatorname{RowNorm}(\bm Y^\xi)^\top\right]_+\\
 &+0.25\left[
 \operatorname{Dist}(\bm U_b,\bm U_n)-0.25\right]_+^{\odot2}
 +0.25\ind[\xi=p]\left[
 \ones_{N_\ell^{\mathrm{OT}}}
 \ones_{N_\ell^{\mathrm{OT}}}^\top
 -\overline{\bm M}_b\overline{\bm M}_n^\top
 \right]_+ .
\end{aligned}}
\label{eq:unified-cost}
\end{equation}
$\overline{\bm M}$ is the stabilized class-wise semantic row normalization.
Only $P$ uses semantic cost.

For $\xi\in\{p,s\}$, define
\begin{equation}
\boxed{
 \bm\Pi^\xi
 =\arg\min_{\bm\Pi\geq0}
 \left\{
 \langle\bm\Pi,\bm C^\xi\rangle_F
 +\varepsilon_\xi\sum_{ij}\Pi_{ij}(\log\Pi_{ij}-1)
 +\mathfrak P_\xi(\bm\Pi)
 \right\},}
\label{eq:unified-transport}
\end{equation}
where
\[
 \mathfrak P_p(\bm\Pi)=
 \begin{cases}
  0,&\bm\Pi\ones=\bm a^p\ \text{and}\
     \bm\Pi^\top\ones=\bm b^p,\\
  +\infty,&\text{otherwise},
 \end{cases}
 \qquad
 \mathfrak P_s(\bm\Pi)=
 \rho_b\KL(\bm\Pi\ones\|\bm a^s)
 +\rho_n\KL(\bm\Pi^\top\ones\|\bm b^s).
\]
The first term is zero when both balanced constraints hold and $+\infty$
otherwise; the latter uses generalized KL.  Thus balanced $P$-OT and
unbalanced $S$-UOT differ only in masses, semantic cost, and marginal
regularization.

For $\bm\varrho^\xi=\bm\Pi^\xi\ones$, the raw expert barycentric teacher is
\[
 \widetilde{\bm Y}^{\xi}
 =\diag(\max(\bm\varrho^\xi,10^{-8}))^{-1}
  \bm\Pi^\xi\bm Y^\xi .
\]
Restoring slice indices, both distillation losses are instances of
\begin{equation}
\boxed{
 \cL_{\xi,\ell}
 =
 \frac{\displaystyle\sum_{m=1}^{M}
  \bm w_{\ell m}^{\xi\top}
  \left[\ones-
  \rowcos(\bm X_{\ell m}^\xi,
  \stopgrad(\widetilde{\bm Y}_{\ell m}^\xi))\right]}
 {\displaystyle\max\!\left(
  \sum_{m=1}^{M}\ones^\top\bm w_{\ell m}^\xi,10^{-8}\right)},
 \quad
 \bm w_{\ell m}^p=\bm a_{\ell m}^p,\quad
 \bm w_{\ell m}^s=\bm\Pi_{\ell m}^s\ones .}
\label{eq:transport-distillation}
\end{equation}
Finally,
$\cL_{P\text{-OT}}=\frac14\sum_{\ell}\cL_{p,\ell}$ and
$\cL_{S\text{-UOT}}=\frac14\sum_{\ell}\cL_{s,\ell}$.  Therefore $P$ weights
valid slices equally, whereas $S$ weights them by received UOT mass.  A
layer-wide received mass at most $10^{-6}$ gives zero $S$ loss.

All masses, error maps, normalized cost copies, transport plans, and
barycentric teachers are detached.  OT gradients reach only the raw
foundation factor tokens.  Expert adapters receive gradients from
reconstruction and decorrelation, never through Sinkhorn.

\section{Structured Training Objective}

The six losses operate on three parts of one pipeline:
\begin{equation}
\boxed{
\begin{aligned}
 \cL(\Theta;e)
 ={}&
 \underbrace{
  \lambda_{\mathrm{seg}}\cL_{\mathrm{seg}}
  +\lambda_{\mathrm{route}}(e)\cL_{\mathrm{route}}
 }_{\text{prompt-conditioned prediction}}\\
 &+\underbrace{
  \lambda_{\mathrm{ae}}\cL_{\mathrm{ae}}
  +\lambda_{\mathrm{ort}}\cL_{\mathrm{ort}}
 }_{\text{complete, non-redundant factorization}}\\
 &+\underbrace{
  \lambda_p(e)\cL_{P\text{-OT}}
  +\lambda_s(e)\cL_{S\text{-UOT}}
 }_{\text{component-specific expert transfer}} .
\end{aligned}}
\label{eq:structured-objective}
\end{equation}
The generic target weights are
$(3,0.2,0.3,10^{-3},0.1,0.1)$.  Router weight reaches $10^{-3}$ linearly over
five epochs.  $P$-OT starts at epoch $2$, $S$-UOT at epoch $3$, and both reach
$0.1$ through a five-epoch linear ramp.

\section{LGE as a Composition of One Student}

The source foreground labels are scar $1$, edema $2$, LV $3$, normal myocardium
$4$, and RV $5$.  The disjoint prompt groups are
\begin{equation}
 \cG^A=((3),(5),(1,2,4)),\qquad
 \cG^R_{\mathrm{split}}=((3),(5),(4),(1),(2)),\qquad
 \cG^R_{\mathrm{merged}}=((3),(5),(4),(1,2)).
 \label{eq:lge-groups}
\end{equation}
Equivalently, let
$\bm Y_{\mathrm{src}}\in\{0,1\}^{V\times5}$ be the foreground one-hot source
label matrix (background/ignored rows are zero), and let
$\bm G_j\in\{0,1\}^{5\times C_j}$ have entries
$[\bm G_j]_{kc}=\ind[k\in G_c^j]$.  Then all prompt targets are obtained by the
single contraction $\bm Y^j=\bm Y_{\mathrm{src}}\bm G_j$.

Let $\mathfrak S_\Theta(\bm x^b;\cG)$ be the pure-student prediction in
Eq.~\eqref{eq:student-master}.  Validation and deployment use the online
composition
\begin{equation}
\boxed{
\begin{aligned}
 \bm L^A
 &=\mathfrak S_\Theta(\bm x^b;\cG^A),\\
 \mathcal R_{\mathrm{online}}^A
 &=\stopgrad\!\left\{
 \RoiBox_{\gamma=1.25}
 \bigl(\ind[\sigm(\bm L^A_{\mathrm{MYO}})\geq0.3]\bigr)
 \right\},\\
 \bm L^R
 &=\mathfrak S_\Theta(
 \Crop(\bm x^b;\mathcal R_{\mathrm{online}}^A);\cG^R).
\end{aligned}}
\label{eq:lge-online}
\end{equation}
The box has minimum side $8$, uses full-image fallback when empty, and is
detached.

Training is a different operator because auxiliary losses require aligned
expert input and target:
\[
 \mathfrak J_\Theta(\bm x^b,\bm x^n,\bm Y,\bm V;\cG)
 =(\bm L,\cL_{\mathrm{seg}},\cL_{\mathrm{aux}}).
\]
Epochs $1$--$10$ evaluate only full-image Pass~A.  Before epoch $11$,
deterministic student predictions create the default fixed ROI cache.  Mixed
training uses
\begin{equation}
 \mathcal R_b^{\mathrm{train}}
 =\operatorname{Jitter}(\mathcal R_b^{\mathrm{cache}}),
 \qquad
 \mathfrak J_\Theta\!\left(
  \Crop_{\mathcal R_b^{\mathrm{train}}}
  (\bm x^b,\bm x^n,\bm Y),\bm V;\cG^R\right),
 \label{eq:lge-cache}
\end{equation}
so the current Pass~A prediction does not choose the training crop.
The image tensors use bilinear interpolation, the target uses nearest-neighbor
interpolation, and the slice-validity vector $\bm V$ is passed through without
spatial cropping.

If a refinement group is present in the full target, its prompt is valid only
when at least $1\%$ of that group lies inside the ROI actually used.  A group
absent from the full slice remains a valid negative.  With
\[
 \cL_{\mathrm{aux}}^j
 =0.2\cL_{\mathrm{ae}}^j+0.3\cL_{\mathrm{ort}}^j
 +\lambda_{\mathrm{route}}(e)\cL_{\mathrm{route}}^j
 +\lambda_p(e)\cL_{P\text{-OT}}^j
 +\lambda_s(e)\cL_{S\text{-UOT}}^j,
\]
the complete two-stage objective is
\begin{equation}
\boxed{
 \cL_{\mathrm{LGE}}(e)=
 \begin{cases}
  \cL_{\mathrm{seg}}^A+\cL_{\mathrm{aux}}^A,
     &e\leq10,\\[1mm]
  \cL_{\mathrm{seg}}^A+\cL_{\mathrm{seg}}^R
  +\dfrac12(\cL_{\mathrm{aux}}^A+\cL_{\mathrm{aux}}^R),
     &e\geq11.
 \end{cases}}
\label{eq:lge-loss}
\end{equation}
Pass~R is not evaluated in the first case.  In mixed training the two forwards
accumulate gradients before one optimizer step.

\subsection{Matrix hard switch}

Let $\bm\iota_R\in\{0,1\}^V$ be the rectangular ROI indicator and
$\bm D_R=\diag(\bm\iota_R)$.  Let $\bm L_\uparrow^R\in\R^{V\times C_R}$ be
refinement logits resized and inserted into the full model grid.  Define
$\bm\Lambda_{AR}\in\{0,1\}^{3\times C_R}$ by
$[\bm\Lambda_{AR}]_{11}=[\bm\Lambda_{AR}]_{22}=1$, and
$\bm q_{AR}=(1,1,0,\ldots,0)^\top$.  Then
\begin{equation}
\boxed{
\begin{aligned}
 \bm L^{\mathrm{out}}
 &=\bm L^A\bm\Lambda_{AR}
   -20\ones_V(\ones_{C_R}-\bm q_{AR})^\top,\\
 \bm L^{\mathrm{final}}
 &=\bm D_R\bm L_\uparrow^R
   +(\bm I-\bm D_R)\bm L^{\mathrm{out}} .
\end{aligned}}
\label{eq:lge-hard-switch}
\end{equation}
Thus every refinement class comes from Pass~R inside the rectangular ROI;
outside it LV/RV come from Pass~A and all other foreground logits are $-20$.
There is no soft blending.  A zero background column is prepended, logits are
restored through the aligned nnUNet geometry, and argmax is taken only after
logit resampling.  The discrete result is then inserted into the saved crop and
inverse-transposed to raw NIfTI geometry.

\section{Training--Deployment Boundary}

The deployed trainable subset is
\[
 \Theta_{\mathrm{deploy}}
 =\{\bm W_{\ell p}^b,\bm W_{\ell s}^b\}_{\ell\in\cI}
  \cup\Theta_{\mathrm{router}}.
\]
It loads the frozen BiomedParse model but not nnUNet, expert autoencoders,
expert preprocessing tensors, mass/cost builders, or Sinkhorn solvers.  LGE
deployment invokes the same pure-student map twice via Eq.~\eqref{eq:lge-online}.

\appendix
\section{Implementation-exact Details}

\paragraph{Inputs and normalization.}
Generic 2.5D BiomedParse input uses
$(x_{z-1},x_z,x_{z+1})$ with boundary replication.  Its CT transform uses
level $40$, width $400$:
$x_i^b=255[\clip(x_i,-160,240)+160]/400$.
LGE uses $Z=1$ and $(x_z,x_z,x_z)$.  It reuses the frozen nnUNet crop/resampling
geometry but computes first/99th percentiles after removing values at or below
the volume's $0.1$th percentile, clips to that range, maps to $[0,255]$, and
stores uint8.  Degenerate percentile ranges use width one.

\paragraph{ROI geometry and augmentation.}
The myocardium threshold is $0.3$; the smallest axis-aligned box is expanded by
$1.25$, shifted/clipped to image bounds, and falls back to the full image when
empty.  Images use bilinear and targets nearest-neighbor resize to $256^2$.
The same training ROI is applied to nnUNet input, BiomedParse input, and target.
Mixed-training cache jitter shifts the center by at most $3\%$ of ROI extent and
enlarges it by a factor in $[1,1.15]$.  Passes receive independent transforms:
rotation $\pm10^\circ$, scale $[0.95,1.05]$, translation $\pm5\%$,
horizontal/vertical flip probabilities $0.5/0.2$, and intensity perturbation
probability $0.8$.  Contrast is $[0.8,1.2]$, shift is $[-0.1,0.1]$ times slice
standard deviation, noise scale is $[0,0.04]$ times that deviation, and
BiomedParse gamma is $[0.75,1.35]$.

\paragraph{Transport and optimization.}
OT uses float32 log-space Sinkhorn for $30$ iterations with
$\varepsilon_p=\varepsilon_s=0.1$, $\rho_b=1.0$, and $\rho_n=0.2$.
AdamW uses learning rate $10^{-4}$, weight decay $10^{-4}$, optional AMP, and
global gradient clipping at $5$.  Generic training keeps a constant learning
rate; LGE uses epoch-wise cosine decay to $0.05$ of its initial value.

\paragraph{Current LGE instantiation.}
At $256^2$, feature sizes are $(64,32,16,8)^2$, nnUNet channels are
$(128,256,512,512)$, BiomedParse channels are $(192,384,768,1536)$, and OT
grids are $(32,32,16,8)^2$.  All trainable modules contain
$10{,}140{,}994$ parameters; the deployment subset contains $6{,}432{,}130$.

\end{document}
